% !TEX root = ../../main.tex
\subsection{Exercise 3}
\task{Requirement}
Given the sequence of word addresses
\texttt{0, 4, 1, 5, 65, 1, 67, 46, 1, 70, 2, 0}. The system has a 256-byte cache,
4-word blocks and the access unit is one byte. Determine the number of HITs and
MISSes for the following cache configurations:
(a) Direct mapped, (b) 2-way set associative, (c) Fully associative.

\task{Solution}
\begin{itemize}[itemsep=2pt]
  \item Block size $= 4$ words $= 16$ bytes; the cache has $256/16 = 16$ blocks.
  \item Block ID $= \lfloor \text{word address} / 4 \rfloor$:
        \texttt{0, 1, 0, 1, 16, 0, 16, 11, 0, 17, 0, 0}.
  \item Set index $=$ Block ID $\bmod$ number of sets; LRU replacement is used.
\end{itemize}

\begin{table}[H]
  \centering
  \small
  \begin{tabular}{lcccccccccccc}
    \toprule
    \textbf{Word address} & 0 & 4 & 1 & 5 & 65 & 1 & 67 & 46 & 1 & 70 & 2 & 0 \\
    \textbf{Block ID}     & 0 & 1 & 0 & 1 & 16 & 0 & 16 & 11 & 0 & 17 & 0 & 0 \\
    \midrule
    (a) DM, 16 sets (index)  & 0 & 1 & 0 & 1 & 0 & 0 & 0 & 11 & 0 & 1 & 0 & 0 \\
    \quad Result             & M & M & H & H & M & M & M & M  & M & M & H & H \\
    (b) 2-way, 8 sets (index)& 0 & 1 & 0 & 1 & 0 & 0 & 0 & 3  & 0 & 1 & 0 & 0 \\
    \quad Result             & M & M & H & H & M & H & H & M  & H & M & H & H \\
    (c) Fully associative    & M & M & H & H & M & H & H & M  & H & M & H & H \\
    \bottomrule
  \end{tabular}
  \caption{Exercise 3 -- HIT (H) / MISS (M) of each access}
\end{table}

\begin{itemize}[itemsep=2pt]
  \item (a) Direct mapped: blocks 0 and 16 both map to index 0 and keep evicting
        each other $\Rightarrow$ \textbf{4 HITs, 8 MISSes}.
  \item (b) 2-way set associative: blocks 0 and 16 fit in the two ways of set 0
        $\Rightarrow$ \textbf{7 HITs, 5 MISSes}.
  \item (c) Fully associative: only the 5 compulsory misses (blocks 0, 1, 16, 11, 17)
        $\Rightarrow$ \textbf{7 HITs, 5 MISSes}.
\end{itemize}

\task{Implementation}
The subroutine \texttt{sim} simulates a set-associative cache with LRU replacement.
Each line stores the block ID (\texttt{-1} means invalid) and the time it was last
used. For every access, the block ID and set index are computed, the ways of the set
are searched; on a miss, the least recently used line (or an invalid line) is replaced.
The routine is called with 1 way, 2 ways and 16 ways (fully associative).
\lstinputlisting[caption={Exercise3.asm}]{Exercise3.asm}

\task{Output}
\screenshot{images/ex3.png}{Exercise 3 -- Output in MARS}
